\documentclass[11pt,a4paper]{article}

\usepackage[utf8]{inputenc}
\usepackage[T1]{fontenc}
\IfFileExists{spanish.ldf}{\usepackage[spanish,es-tabla]{babel}}{}
\usepackage[a4paper, left=1.4cm, right=1.4cm, top=2.5cm, bottom=2.0cm]{geometry}
\usepackage{lmodern}
\IfFileExists{microtype.sty}{\usepackage{microtype}}{}
\usepackage{enumitem}
\usepackage{tabularx}
\usepackage{booktabs}
\usepackage[table]{xcolor}
\usepackage{titlesec}
\usepackage{fancyhdr}

% Definición de tonos sobrios y formales (tipo cuestionario institucional)
\definecolor{grisCabecera}{gray}{0.92}
\definecolor{grisBorde}{gray}{0.35}
\definecolor{textoMuted}{gray}{0.25}

% Configuración de listas compactas con mayor tamaño de fuente
\setlist{nosep, leftmargin=1.35em, topsep=2pt, itemsep=2.5pt}

% Configuración de títulos sobrios y discretos
\titleformat{\section}
  {\normalfont\large\bfseries}
  {}{0em}
  {}
  [{\titlerule[0.7pt]}]

\titlespacing*{\section}{0pt}{8pt}{4pt}

% Encabezados y pies de página formales
\pagestyle{fancy}
\fancyhf{}
\renewcommand{\headrulewidth}{0.4pt}
\renewcommand{\footrulewidth}{0.4pt}
\fancyhead[L]{\small\textbf{UNLP --- Facultad de Ingeniería} \textbar\ Ingeniería de Software 2026 \textbar\ Grupo 7 \textbar\ \textbf{Softech}}
\fancyfoot[L]{\footnotesize Guía de mesa para interlocutores y moderación (relevamiento inicial)}
\fancyfoot[R]{\small Hoja \thepage\ de 2}

\begin{document}

% ENCABEZADO INSTITUCIONAL
\noindent
\begin{tabularx}{\textwidth}{@{}X r@{}}
  \textbf{\large PROTOCOLO Y GUÍA DE RELEVAMIENTO EN MESA} & \textbf{\normalsize ROL: INTERLOCUTORES} \\
  \textit{Empresa: Softech \textbar\ Proyecto: Plataforma de Repartos} & \small Moderación \& Timekeeper \\
\end{tabularx}
\vspace{1pt}
\hrule height 1.1pt
\vspace{4pt}

% TABLA DE METADATOS TIPO CUESTIONARIO DISCRETO (SOLO FILA GRIS SOLICITADA)
\noindent
\begin{tabularx}{\textwidth}{|X|X|X|X|}
\hline
\rowcolor{grisCabecera}
\textbf{Entrevistado:} Martín & \textbf{Fecha:} 04/09/2026 & \textbf{Hora:} 09:20 hs & \textbf{Modalidad:} Presencial \\
\hline
\end{tabularx}

\vspace{5pt}

% APERTURA
\section*{Apertura \& Consentimiento Informado \hfill \normalfont\small\itshape [Interlocutor 1]}
\begin{enumerate}[label=\textbf{\alph*.}]
  \item \textbf{Saludo institucional y encuadre:} Agradecer la presencia y explicar brevemente la meta: conocer la dinámica diaria de su trabajo para relevar los procesos reales.
  \item \textbf{Autorización de grabación:} 
  \textit{``Martín, para trabajar con datos exactos y relevar fielmente el proceso sin perder ningún detalle técnico vamos a grabar el audio de la reunión, ¿estás de acuerdo?''}
  \item \textbf{Identificación protocolar (si autoriza):} Al encender el grabador, cada uno de los 6 integrantes enuncia su \textbf{Nombre y Apellido}, y Martín ratifica oralmente su conformidad para el registro.
  \item \textbf{Plan de contingencia (si no autoriza):} Toma de notas manual exhaustiva por transcriptores y firma de conformidad al pie al finalizar.
\end{enumerate}

% BLOQUE 1
\section*{Bloque 1: Organización, Rol de Martín y Contexto General \hfill \normalfont\small\itshape [Interlocutor 1]}
\begin{itemize}
  \item \textbf{1. Rol y estructura organizativa:}
  \textit{``Martín, para arrancar y situarnos en tu día a día: ¿Cuál es exactamente tu rol dentro de la empresa y cómo está organizada hoy en día?''}
  
  \begin{itemize}[label=$\circ$]
    \item \textit{Bifurcación A1 (Coordinación individual / 'Hombre orquesta'):} 
    ``Al coordinar vos solo todos los pedidos, ¿cómo te organizás con el celular en horas pico para no atrasarte con los mensajes? ¿Sos vos también quien sale a repartir en algún momento?''
    \item \textit{Bifurcación A2 (Equipo gestor / socios / ayudantes):} 
    ``¿Cómo se dividen las tareas o comercios entre ustedes? ¿Manejan una única línea telefónica o varias? ¿Tienen que rendirse cuentas o reportes entre ustedes?''
  \end{itemize}

  \item \textbf{2. Comercios adheridos y rubros:}
  \textit{``¿Con cuántos comercios trabajás actualmente y de qué rubros son (gastronomía, farmacias, tiendas, etc.)?''}

  \item \textbf{3. Equipo de cadetes y medios de transporte:}
  \textit{``¿Cuántos repartidores tenés a cargo y en qué se mueven (bicicleta, moto, auto)?''}\\
  \textit{\small[Repregunta: ¿Son fijos del equipo o trabajan de manera libre/independiente?]}

  \item \textbf{4. Zonas y logística de retiro:}
  \textit{``¿En qué zonas geográficas realizan los repartos? ¿Los paquetes salen directo desde cada local hacia el cliente (punto a punto) o pasan primero por algún lugar donde se juntan o acopian?''}

  \item \textbf{5. Contacto con clientes finales:}
  \textit{``Los clientes que compran en los locales, ¿se comunican alguna vez con vos o con los chicos, o su contacto es únicamente con el comercio?''}
\end{itemize}

% BLOQUE 2
\section*{Bloque 2: Flujo Operativo Actual (AS-IS) \hfill \normalfont\small\itshape [Interlocutor 2]}
\noindent
\textit{\small Estrategia: Formular la pregunta troncal abarcativa para escuchar el circuito natural. Usar el checklist solo para repreguntar vacíos.}

\vspace{3pt}
\noindent
\begin{tabularx}{\textwidth}{|X|}
\hline
\textbf{Pregunta Troncal Abarcativa:}
\textit{``Contanos paso a paso qué pasa desde que un comercio te avisa que tiene un pedido hasta que se entrega al cliente: ¿Qué datos te mandan, cómo se lo asignás al repartidor y cómo finaliza la entrega?''} \\
\hline
\end{tabularx}

\vspace{10pt}

\begin{itemize}[label={[\quad]}]
  \item \textbf{Entrada del pedido:} ¿Por qué canal entra (WhatsApp personal, grupo)? ¿Qué datos envían (dirección exacta, cliente, horario tope, método de pago, tamaño del bulto)? ¿Tenés datos ya agendados o recordados?
  \item \textbf{Asignación al repartidor:} ¿Qué criterio usás para decidir a quién mandar (cercanía en el mapa, orden de llegada a la parada, tipo de vehículo según tamaño)? ¿Cómo le notificás la orden?
  \item \textbf{Monitoreo del estado:} ¿Cómo sabés en qué estado está el pedido a lo largo del día (retirado del local, en viaje, entregado, demora)? ¿Te avisan uno a uno por chat?
  \item \textbf{Cobros y rendiciones:} ¿Ustedes cobran el valor del producto al cliente final o solo el flete/envío? ¿Qué medios usan (efectivo, transferencias)? ¿Con qué frecuencia liquidan a comercios y cadetes?
\end{itemize}

\clearpage

% BLOQUE 3
\section*{Bloque 3: Puntos de Dolor con WhatsApp \& Contingencias \hfill \normalfont\small\itshape [Interlocutor 3]}
\begin{itemize}
  \item \textbf{1. Pérdida de control y errores operativos habituales:}
  \textit{``¿Cuáles son las situaciones más frecuentes donde sentís que se te va de las manos o se pierde el control con WhatsApp?''}
  \begin{itemize}[label=$\circ$]
    \item \textit{Foco de relevamiento:} Mensajes que quedan sin leer o perdidos en horas pico, pedidos duplicados, reclamos cruzados entre comercio y cadete.
  \end{itemize}

  \item \textbf{2. Excepciones e imprevistos en destino:}
  \textit{``¿Qué pasa cuando un repartidor llega y el cliente no atiende el timbre, la dirección es errónea o no tiene cambio? ¿Cómo se resuelve y a quién recurre el cadete en ese momento?''}

  \item \textbf{3. Registros históricos y trazabilidad:}
  \textit{``¿Llevan algún tipo de registro de las entregas que se completaron durante el día o la semana (planillas de cálculo, cuadernos, o no queda registro formal)?''}

  \item \textbf{4. Carga de trabajo manual más crítica:}
  \textit{``En tu día a día, ¿cuál es la tarea manual que más tiempo y energía te consume o que te resulta más desgastante?''}
\end{itemize}

% BLOQUE 4
\section*{Bloque 4: Alcance del Sistema (TO-BE) \& Expectativas \hfill \normalfont\small\itshape [Interlocutor 3]}
\begin{itemize}
  \item \textbf{1. Definición de actores y fronteras del software (Pregunta Clave de Alcance):}
  \textit{``Pensando en el sistema que necesitás: ¿Te imaginás una herramienta de uso exclusivo para vos y la coordinación interna, o la idea es que los comercios y los repartidores también ingresen a cargar o ver estados?''}
  \begin{itemize}[label=$\circ$]
    \item \textit{Repregunta cliente final:} ``¿Y los compradores que reciben el pedido en su casa, necesitarían entrar a ver un estado o seguimiento, o quedan totalmente afuera del sistema?''
    \item \textit{Perfil de usuarios:} ``¿Hay limitaciones técnicas o barreras en el uso del celular por parte de comercios o repartidores?''
  \end{itemize}

  \item \textbf{2. Volumen y proyección de crecimiento:}
  \textit{``¿Cuántos pedidos manejan un día promedio (y un fin de semana)? ¿Cuánto estimás que podría crecer el negocio en los próximos meses?''}

  \item \textbf{3. Funcionalidades irrenunciables y visión de éxito:}
  \textit{``Si el sistema tuviera una sola cosa bien resuelta el primer día, ¿cuál sería la función que para vos no puede faltar bajo ningún punto de vista?''}

  \item \textbf{4. Materiales y artefactos de referencia:}
  \textit{``¿Nos podrías facilitar o reenviar luego capturas de mensajes típicos, listas o planillas que uses hoy para entender el formato real?''}
\end{itemize}

% CIERRE FORMAL
\section*{Cierre Formal, Validación \& Traspaso a Firmas \hfill \normalfont\small\itshape [Cierre Grupal]}
\begin{enumerate}[label=\textbf{\arabic*.}]
  \item \textbf{Síntesis de comprobación:} Resumir en 2 o 3 frases lo comprendido (ej.: \textit{``Entendemos que el núcleo es aliviar la saturación de chats entre comercios y repartidores en hora pico y dar visibilidad clara de qué cadete tiene cada pedido y las cuentas claras al final del día''}) para que Martín asienta y valide.
  \item \textbf{Agradecimiento institucional:} Agradecer sinceramente a Martín por su tiempo y aportes.
  \item \textbf{Paso a firmas (Acta de Relevamiento):} Traspasar formalmente a los transcriptores para la firma de Martín y de los integrantes del equipo en las hojas de notas.
\end{enumerate}

\vspace{8pt}
\hrule
\vspace{5pt}

% RECORDATORIO DISCRETO PARA MODERACIÓN
\noindent
\begin{tabularx}{\textwidth}{@{}l X@{}}
\small\textbf{Pautas de Moderación:} & 
\small
$\bullet$ Tono profesional y empático; concentrarse 100\% en procesos de negocio (sin debatir código). \\
& \small $\bullet$ \textbf{Transición ante seña de tiempo:} \textit{``Martín, clarísimo este punto; para cuidar tus tiempos y cumplir con lo previsto, le doy paso a mi compañero con el siguiente bloque...''}.
\end{tabularx}

\end{document}
